% Bài làm buổi 20/10 — Mít
\begin{problem}[8DS-HSBN-001]
$f(0) = 3 \cdot 0 - 2 = -2$
$f(2) = 3 \cdot 2 - 2 = 4$
$f\left(-\dfrac{1}{3}\right) = 3 \cdot \left(-\dfrac{1}{3}\right) - 2 = -1 - 2 = -3$
Vậy $f(0) = -2$, $f(2) = 4$, $f\left(-\dfrac{1}{3}\right) = -3$.
\end{problem}

\begin{problem}[8DS-HSBN-005]
Vì đồ thị song song với đường thẳng $y = -3x + 5$ nên $a = -3$ và $b \neq 5$.
Đồ thị cắt trục hoành tại điểm có hoành độ bằng $2$ nên đi qua điểm $(2; 0)$.
Thay $x = 2$, $y = 0$ vào $y = -3x + b$: $0 = -3 \cdot 2 + b$
$b = 6$ (thoả mãn $b \neq 5$)
Vậy hàm số cần tìm là $y = -3x + 6$.
\end{problem}

\begin{problem}[8DS-PT1-006]
Gọi tuổi con năm nay là $x$ (tuổi), điều kiện $x$ nguyên dương.
Tuổi mẹ năm nay là $3x$ (tuổi).
Sau $12$ năm nữa tuổi con là $x + 12$, tuổi mẹ là $3x + 12$.
Theo đề bài ta có phương trình: $3x + 12 = 2(x + 12)$
$3x + 12 = 2x + 24$
$y = 12$ (thoả mãn điều kiện)
Vậy năm nay con $12$ tuổi, mẹ $36$ tuổi.
\end{problem}

\begin{problem}[8DS-PS-001]
a) Phân thức $\dfrac{3x}{x - 2}$ xác định khi $x - 2 \neq 0$, tức là $x \neq 2$.
b) Phân thức $\dfrac{x + 1}{x^2 - 9}$ xác định khi $x^2 - 9 \neq 0$
$(x - 3)(x + 3) \neq 0$
$x \neq 3$
Vậy ĐKXĐ: a) $x \neq 2$; b) $x \neq 3$.
\end{problem}
